\subsection{Endomorphisms of finite trace}

We recall that a positive endomorphism \(u\) of E is of finite trace if and only if there exists an orthonormal basis \((e_i)_{i \in I}\) of E such that

\[
\sum_ {i \in \mathrm{I}} \langle e _ {i} | u (e _ {i}) \rangle <   + \infty
\]

(EVT, V, p. 48, lemma 3 and p. 49, def. 7). If K = C, the space \(\mathcal{L}_1(\mathrm{E})\) of endomorphisms of finite trace of E is the vector space generated by the set of positive endomorphisms of finite trace (EVT, V, p. 50, def. 8); if K = R, the space \(\mathcal{L}_1(\mathrm{E})\) is defined as the intersection \(\mathcal{L}(\mathrm{E}) \cap \mathcal{L}_1(\mathrm{E}_{(\mathbf{C})})\) (EVT, V, p. 50).

If \(u\in \mathcal{L}_1(\mathrm{E})\) , then the series

\[
\sum_ {i \in \mathrm{I}} \langle e _ {i} | u (e _ {i}) \rangle
\]

converges for every orthonormal basis \((e_{i})_{i\in\mathrm{I}}\) of E and its sum is independent of the orthonormal basis; this sum is called the trace \(\operatorname{Tr}(u)\) of u (EVT, V, p. 50). If K = R, we have \(\operatorname{Tr}(u) = \operatorname{Tr}(u_{(\mathbf{C})})\) .

Let \(u \in \mathcal{L}_1(\mathrm{E})\) . We have \(u^* \in \mathcal{L}_1(\mathrm{E})\) and \(\operatorname{Tr}(u^*) = \overline{\operatorname{Tr}(u)}\) (loc. cit.).

PROPOSITION 12. --- Let \(B = (e_i)_{i \in I}\) be an orthonormal basis of E. Let \(u \in \mathcal{D}_B(E)\) and denote by \(\lambda = (\lambda_i)_{i \in I}\) the family of its eigenvalues. The endomorphism \(u\) is of finite trace if and only if the family \(\lambda\) is summable in K. We then have \(\operatorname{Tr}(u) = \sum \lambda_i\) .

Replacing \(u\) by \(u_{(\mathbf{C})}\) if necessary, we may assume that \(K = \mathbf{C}\) .

Suppose first that \(u\) is of finite trace. According to EVT, V, p. 50 and p. 49, formula (25), the family \((\langle e_i | u(e_i) \rangle)_{i \in I} = (\lambda_i)_{i \in I}\) is summable.

Conversely, suppose that the family \(\lambda\) is summable.

Each of the families \((\mathcal{R}(\lambda_{i})^{+})\) , \((\mathcal{R}(\lambda_{i})^{-})\) , \((\mathcal{I}(\lambda_{i})^{+})\) , \((\mathcal{I}(\lambda_{i})^{-})\) is then summable. The endomorphism u is a linear combination of the elements of \(\mathcal{D}_{\mathrm{B}}(\mathrm{E})\) whose eigenvalues are these families. These elements of \(\mathcal{D}_{\mathrm{B}}(\mathrm{E})\) are positive. Since, by definition, the space of finite-trace endomorphisms is generated by positive finite-trace endomorphisms, we may therefore assume that \(\lambda_{i} \geqslant 0\) for all i. Since \(\langle e_{i} | u(e_{i}) \rangle = \lambda_{i}\), the family \((\langle e_{i} | u(e_{i}) \rangle)_{i \in \mathrm{I}}\) is summable, hence u is of finite trace (EVT, V, p. 48, lemma 3).

Finally, if u has finite trace, then by EVT, V, p. 50, we have

\[
\operatorname{Tr} (u) = \sum_ {i \in \mathrm{I}} \left\langle e _ {i} \mid u \left(e _ {i}\right) \right\rangle = \sum_ {i \in \mathrm{I}} \lambda_ {i}.
\]

COROLLARY 1. --- Let u be an endomorphism of E with finite trace.

\begin{enumerate}
\def\labelenumi{\alph{enumi})}
\item
  The endomorphism u is compact;
\item
  Let \(\mathrm{B} = (e_i)_{i \in \mathrm{I}}\) be an orthonormal basis of the initial space \(\operatorname{Ker}(u)^{\circ}\) of u, \(\mathrm{C} = (f_i)_{i \in \mathrm{I}}\) be an orthonormal family of E and \((\alpha_i)_{i \in \mathrm{I}}\) be a family in \((\mathbf{R}_+^*)^{\mathrm{I}}\) such that \(u(e_i) = \alpha_i f_i\) for all \(i \in I\) (cor. 2 of IV, p. 150). We have
\end{enumerate}

\[
\operatorname{Tr} (u) = \sum_ {i \in I} \alpha_ {i} \langle e _ {i} | f _ {i} \rangle .
\]

To prove \(a\) ), one may assume that \(\mathbf{K} = \mathbf{C}\) (EVT, V, p. 50 and remark 4 of III, p. 2) and that \(u\) is positive (EVT, V, p. 50, def. 8).

By EVT, V, p. 56, cor. 1, there exists an orthonormal basis \(B = (e_i)_{i \in I}\) of E such that \(u\) is diagonal in the basis B and moreover the family \(\lambda\) of the eigenvalues of \(u\) belongs to \(\mathbf{R}_+^I\) and is summable. The family \(\lambda\) therefore belongs to \(\mathcal{C}_0(\mathrm{I};\mathrm{K})\) (TG, III, p. 38, prop. 1), and consequently the endomorphism \(u\) is compact (prop. 2 of IV, p. 148).

Let us prove \(b\) ). Let \((e_i)_{i \in \mathbf{J}}\) be an orthonormal basis of E extending the family B, so that \(u(e_i) = 0\) if \(i \in \mathbf{J} - \mathbf{I}\) . We obtain

\[
\operatorname{Tr} (u) = \sum_ {i \in \mathrm{J}} \langle e _ {i} \mid u (e _ {i}) \rangle = \sum_ {i \in \mathrm{I}} \langle e _ {i} \mid u (e _ {i}) \rangle = \sum_ {i \in \mathrm{I}} \alpha_ {i} \langle e _ {i} \mid f _ {i} \rangle .
\]

COROLLARY 2. --- Let F be a Hilbert space. Let \(u \in \mathcal{L}(\mathrm{E};\mathrm{F})\) be a Hilbert--Schmidt operator. Then u is a compact linear operator.

Let \((j,|u|)\) be the polar decomposition of u (I, p. 140, def. 4). By definition (EVT, V, p. 50, def. 9) the endomorphism \(u^{*}u\) has finite trace, hence is compact (corollary 1, a)); the same is true of \(|u|=\sqrt{u^{*}u}\) (prop. 6 of III, p. 91, b)) and of \(u=j|u|\) (prop. 3 of III, p. 5).

COROLLARY 3. --- Suppose that K = C. Let u be a positive compact endomorphism of E. The endomorphism u is of finite trace if and only if the decreasing family of its eigenvalues \((\lambda_{n}(u))_{n\in\mathrm{I}_{\mathrm{E}}}\) is summable; the trace of u is then the sum of this family.

By th. 1 of IV, p. 149, this follows from the preceding proposition and the definition of the sequence \((\lambda_n(u))_{n\in \mathrm{I_E}}\) (no. 3 of IV, p. 151), taking into account formula (28) of EVT, V, p. 49.

COROLLARY 4. --- Let F be a Hilbert space. Let \(u \in \mathcal{L}(\mathrm{E};\mathrm{F})\) be a compact linear map. Let \(\mathrm{B} = (e_i)_{i \in \mathrm{I}}\) be an orthonormal basis of the initial space \(\operatorname{Ker}(u)^{\circ}\) of \(u\) , \(\mathrm{C} = (f_i)_{i \in \mathrm{I}}\) be an orthonormal family of F and \((\alpha_{i})_{i\in\mathrm{I}}\) be a family in \((\mathbf{R}_{+}^{*})^{\mathrm{I}}\) such that \(u(e_{i})=\alpha_{i}f_{i}\) for all \(i\in I\) (cor. 2 of IV, p. 150).

The endomorphism \(|u|\) of E has finite trace if and only if the family \((\alpha_{i})_{i\in \mathrm{I}}\) is summable. We then have \(u\in \mathcal{L}_2(\mathrm{E};\mathrm{F})\) and

\[
\operatorname{Tr} (| u |) = \sum_ {i \in \mathrm{I}} \alpha_ {i}, \quad \| u \| _ {2} ^ {2} = \operatorname{Tr} (u ^ {*} u) = \sum_ {i \in \mathrm{I}} \alpha_ {i} ^ {2},\tag{8}
\]

and in particular \(\|u\|_{2}\leqslant\mathrm{Tr}(|u|)\) .

The family of nonzero eigenvalues of \(|u|\) is \((\alpha_{i})_{i\in \mathrm{I}}\) , hence \(|u|\) has finite trace if and only if the family \((\alpha_{i})_{i\in \mathrm{I}}\) is summable (prop. 12). If that is the case, the family \((\alpha_{i}^{2})\) of nonzero eigenvalues of \(u^{*}u = |u|^{2}\) is summable, and formulas (8) result from loc. cit.

Lemma 7. --- Let \(\lambda = (\lambda_i)_{i \in \mathrm{I}}\) be a family of complex numbers. For every \(t \in \mathbf{C}^*\) , let \(n_t\) be the cardinality of the set of \(i \in \mathrm{I}\) such that \(\lambda_i = t\) . The family \((\lambda_i)_{i \in \mathrm{I}}\) is summable if and only if \(n_t\) is finite for all \(t \in \mathbf{C}^*\) and if the family \((n_t t)_{t \in \mathbf{C}^*}\) is summable. In this case, the sums of these two families are equal.

Suppose that the family \((\lambda_{i})_{i\in\mathrm{I}}\) is summable. For every \(t\in\mathbf{C}\) , let \(I_{t}\) be the set of \(i\in I\) such that \(\lambda_{i}=t\) . By TG, III, p. 39, th. 2, applied to the partition of I by the sets \(I_{t}\) , the set \(I_{t}\) is finite for all \(t\in\mathbf{C}^{*}\) , and moreover the family \((n_{t}t)_{t\in\mathbf{C}^{*}}\) is summable with sum equal to that of the family \((\lambda_{i})_{i\in\mathrm{I}}\) .

Conversely, suppose that \(n_{t}\) is finite for all \(t \in \mathbf{C}^{*}\) and that the family \((n_{t}t)_{t \in \mathbf{C}^{*}}\) is summable. Let J be a finite subset of I and \(\Lambda\) the set of \(\lambda_{i}\) for \(i \in J\) . We have

\[
\left| \sum_ {i \in J} \lambda_ {i} \right| \leqslant \sum_ {t \in \Lambda - \{0 \}} n _ {t} | t | \leqslant \sum_ {t \in \mathbf {C} ^ {*}} n _ {t} | t |,
\]

therefore the family \((\lambda_{i})_{i\in\mathrm{I}}\) is summable (TG, VII, p. 17, corollary).

PROPOSITION 13. --- Suppose that K = C. Let u be a compact normal endomorphism of E. For \(t \in \mathrm{Sp}_{\mathrm{s}}(u)\) , let \(n_{t} \geqslant 1\) be the spectral multiplicity of the eigenvalue t of u. For u to have finite trace, it is necessary and sufficient that the family \((n_{t}t)_{t \in \mathrm{Sp}_{\mathrm{s}}(u)}\) be summable. The trace of u is then the sum of this family.

Let \(\mathrm{B} = (e_i)_{i \in \mathrm{I}}\) be an orthonormal basis of E such that u is diagonal in the basis B (th. 1 of IV, p. 149), and let \(\lambda = (\lambda_i)_{i \in \mathrm{I}}\) be the family of its eigenvalues. Since, for \(t \in \mathrm{Sp}_\mathrm{s}(u)\) , the multiplicity \(n_t\) is equal to the number of elements \(i \in I\) such that \(\lambda_i = t\), and since the nonzero elements of \(\lambda\) belong to \(\mathrm{Sp}_{\mathrm{s}}(u)\) , the assertion then follows from prop. 12 and the preceding lemma.

